%%================================================
%% Chapter 1
%%================================================
\chapter{บทนำ}
\label{chapter1}

\section{ที่มาและความสำคัญของปัญหา}
\label{sec:ch1_significance}

ประเทศไทยกำลังเผชิญการเปลี่ยนแปลงโครงสร้างประชากรเข้าสู่สังคมสูงวัยอย่างต่อเนื่อง ซึ่งทำให้ความต้องการระบบสนับสนุนการดูแลผู้สูงอายุและผู้ที่มีข้อจำกัดด้านการเคลื่อนไหวเพิ่มขึ้นตามไปด้วย \cite{SituationThaiElderly2022,WHO2021DecadeAgeing} ขณะเดียวกันสถานการณ์คนพิการในประเทศไทย ปี พ.ศ. 2567 พบว่ามีประชากรผู้พิการกว่า 2.2 ล้านคน โดยผู้พิการทางการเคลื่อนไหวเป็นกลุ่มที่ต้องอาศัยอุปกรณ์ช่วยเคลื่อนที่และการดูแลจากบุคลากรหรือญาติอย่างใกล้ชิด \cite{MSDHS_Disability_2024} แนวคิดการดูแลแบบ Ambient Assisted Living จึงถูกนำมาใช้เป็นกรอบในการออกแบบสภาพแวดล้อมที่ช่วยลดภาระการดูแลและสนับสนุนการใช้ชีวิตอย่างปลอดภัยมากขึ้น \cite{Rashidi2013AALSurvey}

ผู้ใช้เก้าอี้รถเข็นและผู้สูงอายุในสถานดูแล เช่น บ้านพักคนชราวาสนะเวศม์ ซึ่งเป็นบริบทศึกษาของโครงงานนี้ ต้องพึ่งพาผู้ดูแลในการติดตามตำแหน่ง การเคลื่อนที่ และสภาพร่างกายเบื้องต้น ปัญหาดังกล่าวไม่ได้เกิดจากการขาดการดูแลเพียงอย่างเดียว แต่เกิดจากข้อมูลที่กระจายอยู่คนละส่วน เช่น ตำแหน่งผู้พักอาศัย การเคลื่อนไหวบนเก้าอี้รถเข็น สัญญาณชีพจากอุปกรณ์สวมใส่ และบันทึกเวรของผู้ดูแล เมื่อข้อมูลเหล่านี้ไม่ถูกรวมไว้ในระบบเดียว ผู้ดูแลจึงต้องตรวจสอบหลายช่องทางและอาจตอบสนองต่อเหตุการณ์ผิดปกติได้ช้ากว่าที่ควร

\begin{figure}[htbp]
    \centering
    \includegraphics[width=0.75\textwidth]{chapter1/ch1-fig04.jpg}
    \captionsetup{justification=centering}
    \caption{บริบทการดูแลผู้สูงอายุและผู้พิการที่ใช้เก้าอี้รถเข็นในสถานดูแลของประเทศไทย}
    \label{fig:ch1-fig01-thai-context}
\end{figure}

จากการลงพื้นที่สำรวจอาคารและพื้นที่บริการหลัก พบว่าผู้พักอาศัยและจุดให้บริการกระจายอยู่ในหลายอาคาร การติดตามด้วยสายตาหรือการเดินตรวจเพียงอย่างเดียวจึงไม่เพียงพอในทุกช่วงเวลา โดยเฉพาะผู้ใช้เก้าอี้รถเข็นที่อาจเคลื่อนที่ออกจากโซนเดิมหรือหยุดนิ่งเป็นเวลานานโดยไม่มีผู้ดูแลสังเกตเห็นทันที สถานการณ์ลักษณะนี้อาจเพิ่มความเสี่ยงต่ออุบัติเหตุ เช่น การหกล้ม หรือการไม่ได้รับความช่วยเหลือเมื่อเกิดเหตุผิดปกติ \cite{Sterling2016FallsInCareHomes}

\begin{figure}[htbp]
    \centering
    \includegraphics[width=0.75\textwidth]{chapter1/ch1-fig02.jpg}
    \captionsetup{justification=centering}
    \caption{อาคารและพื้นที่บริการหลักของบ้านพักคนชราวาสนะเวศม์ที่ใช้เป็นบริบทศึกษา}
    \label{fig:ch1-fig02-facility-context}
\end{figure}

นอกจากนี้ การทำงานของผู้ดูแลยังมีข้อจำกัดจากการบันทึกข้อมูลแบบแอนะล็อก เช่น กระดานหรือกระดาษ ซึ่งสะดวกต่อการใช้งานเฉพาะหน้าแต่ยากต่อการค้นย้อนหลังและการวิเคราะห์แนวโน้ม ข้อมูลที่แยกจากกันยังทำให้ผู้ดูแลต้องเป็นผู้รวบรวมและตีความเองเป็นส่วนใหญ่ หากมีการแจ้งเตือนจำนวนมากโดยไม่มีบริบทประกอบก็อาจนำไปสู่ความล้าจากการแจ้งเตือน (alert fatigue) และทำให้การจัดลำดับความสำคัญของงานยากขึ้น \cite{Gates2018AlertFatigue}

\begin{figure}[htbp]
    \centering
    \includegraphics[width=0.75\textwidth]{chapter1/ch1-fig03.jpg}
    \captionsetup{justification=centering}
    \caption{การใช้กระดานบันทึกข้อมูลแบบแอนะล็อกที่พบในหน้างานจริง}
    \label{fig:ch1-fig03-analog-records}
\end{figure}

แม้ปัจจุบันจะมีระบบอัตโนมัติในบ้านและระบบระบุตำแหน่งภายในอาคารหลายรูปแบบ \cite{Home_Automation_Integration_2023, Indoor_GPS_Limitations, UWB_Accuracy_Spec} แต่ระบบเหล่านี้มักแยกเป็นส่วนเฉพาะทาง เช่น ระบบติดตามตำแหน่ง ระบบสัญญาณเตือน หรือระบบบันทึกข้อมูลสุขภาพ ทำให้ยังไม่ตอบโจทย์การดูแลที่ต้องเห็นตำแหน่ง การเคลื่อนไหว สัญญาณชีพ และสถานะงานของผู้ดูแลในภาพรวมเดียว โครงงานนี้จึงเสนอการพัฒนาระบบต้นแบบ WheelSense เพื่อรวมข้อมูลจากเซ็นเซอร์บนเก้าอี้รถเข็น อุปกรณ์สวมใส่ และ mobile gateway เข้าสู่ backend และ dashboard ส่วนกลาง พร้อมใช้ AI/MCP เป็นเครื่องมือช่วยค้นข้อมูล วิเคราะห์บริบท และเสนอแนวทางการดำเนินการให้ผู้ดูแลตรวจสอบก่อนใช้งานจริง โดยระบบยังคงอยู่ในระดับต้นแบบเพื่อสนับสนุนการเฝ้าระวังเบื้องต้น ไม่ใช่ระบบวินิจฉัยหรือเครื่องมือแพทย์ที่ได้รับการรับรอง

\section{วัตถุประสงค์ของโครงงาน}
\label{sec:ch1_objectives}

\begin{enumerate}
    \item เพื่อพัฒนาระบบต้นแบบ WheelSense สำหรับสนับสนุนการดูแลผู้ใช้เก้าอี้รถเข็นและผู้สูงอายุในสภาพแวดล้อมแบบ smart environment
    \item เพื่อเก็บข้อมูลจาก wheelchair sensor อุปกรณ์สวมใส่ และ mobile gateway ได้แก่ ข้อมูลการเคลื่อนไหว สัญญาณชีพเบื้องต้น ตำแหน่งภายในอาคาร และข้อมูลการก้าวเดินของผู้ใช้ที่ไม่ได้จับคู่กับเก้าอี้รถเข็น
    \item เพื่อพัฒนา backend และฐานข้อมูลสำหรับรับ ส่ง จัดเก็บ และเรียกดูข้อมูล telemetry จากอุปกรณ์ภาคสนามและแอปพลิเคชันบนอุปกรณ์เคลื่อนที่
    \item เพื่อแสดงผลข้อมูลผู้พักอาศัยและสถานะระบบผ่าน dashboard ตามบทบาทของผู้ใช้งาน เช่น ผู้ดูแล หัวหน้าพยาบาล ผู้สังเกตการณ์ และผู้ป่วย
    \item เพื่อบูรณาการ AI/MCP สำหรับช่วยค้นข้อมูล วิเคราะห์บริบท และเสนอ action ภายใต้ขั้นตอนให้ผู้ดูแลตรวจสอบและยืนยันก่อนดำเนินการ
    \item เพื่อประเมินระบบต้นแบบในด้านการทำงานของเซ็นเซอร์ การระบุตำแหน่งภายในอาคาร ความเสถียรของ mobile gateway/server การใช้งาน dashboard และประสิทธิภาพของ AI pipeline
\end{enumerate}

\section{ขอบเขตของโครงงาน}
\label{sec:ch1_scope}

\begin{enumerate}
    \item ด้านฮาร์ดแวร์ ครอบคลุมการใช้ M5StickC Plus2 เป็นเซ็นเซอร์บนเก้าอี้รถเข็น การใช้ Polar Verity Sense สำหรับอ่านสัญญาณชีพผ่าน BLE และการใช้ Raspberry Pi 5 RAM 8GB เป็น local/edge server สำหรับบริการหลักของระบบ
    \item ด้านการสื่อสาร ครอบคลุมการเชื่อมต่อ BLE ระหว่างอุปกรณ์ภาคสนามกับ mobile gateway และการส่งข้อมูลเข้าสู่ระบบผ่าน MQTT หรือ REST API ตามบทบาทของข้อมูลที่รับส่ง
    \item ด้านซอฟต์แวร์ ครอบคลุม Flutter mobile gateway, FastAPI backend, PostgreSQL 16 พร้อม JSONB, Next.js dashboard, Docker และ MQTT broker สำหรับรับ ส่ง จัดเก็บ และแสดงผลข้อมูล
    \item ด้านการระบุตำแหน่งภายในอาคาร ครอบคลุมการใช้ค่า BLE RSSI จากโหนดอ้างอิง ร่วมกับแนวคิด fingerprinting, KNN และ max-RSSI fallback ในระดับโซนหรือห้อง
    \item ด้าน AI ครอบคลุม pipeline แบบ 5 ชั้น ได้แก่ Intent Router, Context Assembler, Behavioral State, Constrained LLM และ Execution Gate รวมถึงแนวคิด RAG และ MCP tool execution ภายใต้ข้อจำกัดด้านสิทธิ์และการยืนยันก่อนดำเนินการ
    \item ด้านการประเมินผล ครอบคลุมการทดสอบระดับระบบต้นแบบ เช่น sensor test, KNN localization, mobile gateway reliability, dashboard/HMI, Inter-Rater Reliability, cosine similarity, latency และ SUS Score ในกรณีที่มีข้อมูลแบบประเมินรองรับ
    \item โครงงานนี้ไม่มุ่งหมายให้เป็นระบบวินิจฉัยทางการแพทย์ ระบบเฝ้าระวังที่ผ่านการรับรองเป็น medical device หรือการทดลองทางคลินิกกับผู้ป่วยจริง ผลการทดสอบจึงใช้เพื่อประเมินความเป็นไปได้ของระบบต้นแบบเท่านั้น
\end{enumerate}

\section{ประโยชน์ที่คาดว่าจะได้รับ}
\label{sec:ch1_benefits}

\textbf{ประโยชน์ต่อสังคม}
\begin{enumerate}
    \item ช่วยลดภาระการตรวจสอบซ้ำของผู้ดูแล โดยรวมข้อมูลตำแหน่ง การเคลื่อนไหว สัญญาณชีพเบื้องต้น และสถานะอุปกรณ์ไว้ในระบบส่วนกลาง
    \item ช่วยให้ผู้ดูแลเห็นข้อมูลผู้พักอาศัยผ่าน dashboard ตามบทบาท ทำให้สามารถติดตามสถานะและค้นย้อนหลังได้สะดวกกว่าการใช้บันทึกแบบแยกส่วน
    \item ช่วยใช้ข้อมูลจาก sensor เป็นข้อมูลประกอบการประเมินสถานะเบื้องต้น เช่น การหยุดนิ่งนานผิดปกติ การเคลื่อนที่ออกจากโซน หรือค่าสัญญาณชีพที่ควรตรวจสอบเพิ่มเติม
    \item ช่วยเป็นแนวทางสำหรับการพัฒนาระบบดูแลผู้สูงอายุแบบ smart environment ที่สามารถนำไปต่อยอดกับ workflow การแจ้งเตือนและการจัดการงานของผู้ดูแลในอนาคต
\end{enumerate}

\textbf{องค์ความรู้ที่ได้รับ}
\begin{enumerate}
    \item ได้องค์ความรู้ด้านการออกแบบระบบ IoT ที่เชื่อม wheelchair sensor, wearable sensor, mobile gateway, backend และ dashboard ในระบบต้นแบบเดียว
    \item ได้องค์ความรู้ด้านการใช้ BLE RSSI และ KNN สำหรับการระบุตำแหน่งภายในอาคารระดับโซนหรือห้อง
    \item ได้องค์ความรู้ด้านการออกแบบ AI/MCP pipeline ที่แยกการค้นข้อมูล การจัดบริบท การใช้ LLM และการตรวจสอบสิทธิ์ก่อนสั่งงาน
\end{enumerate}

\clearpage
\section{แผนการดำเนินงาน}
\label{sec:ch1_timeline}

การดำเนินงานวิจัยครอบคลุมระยะเวลา 12 เดือน โดยแบ่งเป็นช่วงการวิเคราะห์ปัญหา การศึกษาเอกสารที่เกี่ยวข้อง การออกแบบและสร้างระบบต้นแบบ การบูรณาการและทดสอบ ตลอดจนการวิเคราะห์ผลและจัดทำเล่มวิทยานิพนธ์ ดังรายละเอียดในตารางที่~\ref{tab:chapter1_timeline}

\begin{table}[H]
    \centering
    \caption{แผนการดำเนินงานของโครงงาน}
    \label{tab:chapter1_timeline}
    \small
    \renewcommand{\arraystretch}{1.35}
    \setlength{\tabcolsep}{2.2pt}
    \begin{tabular}{|L{5.8cm}|C{0.45cm}|C{0.45cm}|C{0.45cm}|C{0.45cm}|C{0.45cm}|C{0.45cm}|C{0.45cm}|C{0.45cm}|C{0.45cm}|C{0.45cm}|C{0.45cm}|C{0.45cm}|}
        \hline
        \multirow{2}{*}{\centering \textbf{กิจกรรม}} & \multicolumn{12}{c|}{\textbf{เดือนที่}} \\ \cline{2-13}
        & \textbf{1} & \textbf{2} & \textbf{3} & \textbf{4} & \textbf{5} & \textbf{6} & \textbf{7} & \textbf{8} & \textbf{9} & \textbf{10} & \textbf{11} & \textbf{12} \\ \hline
        1. วิเคราะห์ปัญหาและเก็บบริบทหน้างาน & \(\to\) & \(\to\) & & & & & & & & & & \\ \hline
        2. ศึกษาวรรณกรรมและทฤษฎีที่เกี่ยวข้อง & \(\to\) & \(\to\) & \(\to\) & & & & & & & & & \\ \hline
        3. ออกแบบสถาปัตยกรรมระบบและ data flow & & \(\to\) & \(\to\) & \(\to\) & & & & & & & & \\ \hline
        4. พัฒนา wheelchair sensor และ wearable gateway & & & \(\to\) & \(\to\) & \(\to\) & \(\to\) & & & & & & \\ \hline
        5. พัฒนาระบบระบุตำแหน่งภายในอาคาร & & & & \(\to\) & \(\to\) & \(\to\) & \(\to\) & & & & & \\ \hline
        6. พัฒนา backend, database และ MQTT service & & & & \(\to\) & \(\to\) & \(\to\) & \(\to\) & \(\to\) & & & & \\ \hline
        7. พัฒนา dashboard และ HMI ตามบทบาทผู้ใช้ & & & & & \(\to\) & \(\to\) & \(\to\) & \(\to\) & \(\to\) & & & \\ \hline
        8. บูรณาการ AI/MCP pipeline และ execution gate & & & & & & \(\to\) & \(\to\) & \(\to\) & \(\to\) & & & \\ \hline
        9. ทดสอบ sensor, localization, gateway และ server & & & & & & & & \(\to\) & \(\to\) & \(\to\) & & \\ \hline
        10. ประเมิน dashboard, AI และสรุปผลการทดลอง & & & & & & & & & \(\to\) & \(\to\) & \(\to\) & \\ \hline
        11. จัดทำและปรับปรุงปริญญานิพนธ์ & & & & & & & & & & \(\to\) & \(\to\) & \(\to\) \\ \hline
    \end{tabular}
    \\
    \vspace{0.5em}
    \footnotesize \textbf{หมายเหตุ:} \(\to\) หมายถึงช่วงเวลาที่ดำเนินกิจกรรม
\end{table}
